% 06_soil_exchange.tex - GWSoilExchange.csv and plot_soilzone.py
\section{Look at the exchange between soil and aquifer}

\step{Read \texttt{GWSoilExchange.csv} and plot it}
\cmd{GWSoilExchange.csv} has one line per HRU and day. Show the first columns of the first lines:
\term{t3_soilex_head.txt}
\why the columns are: the base of the soil zone, the soil water level and the water table (means over the HRU's
columns), the saturated fraction (the share of the HRU where the water table reaches the soil), and the exchange in both
directions. The script \cmd{plot\_soilzone.py} draws them for the HRUs you give (here 6 and 1):
\term{t3_plot_soilzone.txt}
\fig{t3_soilzone.png}{\cmd{soilzone.png} for the valley HRUs 6 and 1. Left: the water table rises above the base of the
soil zone every winter, and the soil water level follows it. Right: the exchange (aquifer to soil is positive) and the
saturated fraction (dotted). Bottom: the totals of all HRUs, the rejected recharge and the overflow, and the work of the
daily iteration.}
\begin{tocheck}
\begin{itemize}
\item Once the water table is in the soil, the soil water level follows it (left panels).
\item The exchange changes sign only at real turning points of the seasons, not from one day to the next.
\item The saturated fraction of HRU~6 rises and falls with the seasons (between \V{t3:F6min} and \V{t3:F6max}, mean
  \V{t3:F6mean}). With 25 DEM points per cell, the low spots of the valley floor saturate first.
\item The number of interface updates per day stays small, and the last change of the soil level stays below the
  tolerance (bottom panels).
\end{itemize}
\end{tocheck}
